\honest{Leave a Line of Retreat}{Eliezer Yudkowsky}{2008}

I open with Sun Tzu: when you surround the enemy, leave them an escape route.

At a rationalists' gathering I met a woman who did not believe in cryonics because the soul would not stay with the frozen body. I asked, ``But how do you know that?'' From her face it was ``pretty clear'' she had never asked herself.

I advised her to picture the world without souls, and what she would do in it, before weighing the evidence. Then she could think, ``Well, if there are no souls, I can just sign up for cryonics.'' This is my principle: leave yourself a line of retreat. A layoff is less frightening once you have counted your savings and checked the job market, and only then can you judge the risk ``fairly.'' \nb{Seneca gave a friend facing a lawsuit the same exercise: ``assume that what you fear may happen will certainly happen in any event'' and ``estimate the amount of your fear.'' The use the post adds is to lower the fear before judging the probability. It names no precedent.}

The mechanism is ``the hope'': imagining a possibility takes less courage than admitting it is likely, and makes the admission easier. Imagining a belief is not conceding it. ``The thought you cannot think controls you more than thoughts you speak aloud.'' \nb{That sentence is Yudkowsky's own, from ``Twelve Virtues of Rationality.''}

You must at least admit which ideas scare you. ``Does it help if I say that I have occasion to use this technique myself?'' I describe no occasion.

Think of all the people who say that without God, morality is impossible. If they pictured their real reaction, they would see that ``they wouldn't go around slaughtering babies.'' A footnote says the topic came up in the conversation, so I am ``not offering a strawman.'' \nb{As philosophers make it, the claim is about whether anything would be objectively right or wrong without God, not about how anyone would behave. Some atheists, Nietzsche among them, have accepted a version of it.}

People get over things. ``Newly minted quadriplegics are not as sad, six months later, as they expect to be,'' and ``Quadriplegics deal, and so would you.'' I give no source. \nb{Research supports the direction, and also finds the recovery partial: in two national panels, a long-term disability brought ``moderate to large drops in happiness \ldots\ followed by little adaptation over time'' (Lucas, 2007).}

What is true is already so; owning up to it does not make it worse, and picturing a world you fear does no harm either way. \nb{For the limits of the first claim, see ``You Can Face Reality.'' The second, about picturing rather than learning, does not share them.}

How many religious people would keep their belief if they could accurately picture a world without God? Leaving a line of retreat ``is a powerful technique, but it's not easy.'' I do not say how the conversation ended.
